% Integral bridge for the actual exponent-gluing groups; wrapper-free.
The following construction gives the exact relation to degree-zero cyclic
homology and cyclic cohomology, including the integral torsion and the
original bounded exponent fibres. All group algebras here are over \(\Z\).

\begin{lemma}[Integral augmentation coordinates and their complete fibres]
\label{lem:integral-augmentation}
Let \(G\) be an abelian group written additively. Put
\(A_G=\Z[G]=\bigoplus_{g\in G}\Z u_g\), with
\(u_gu_h=u_{g+h}\) and unit \(u_0\), and put
\(I_G=\ker(\epsilon_G:A_G\to\Z)\), where
\(\epsilon_G(\sum_g n_gu_g)=\sum_g n_g\).
Every sum here has finite support. The maps
\[
 \kappa_G:G\longrightarrow I_G/I_G^2,
 \qquad g\longmapsto (u_g-u_0)+I_G^2,
\]
\[
 \theta_G:I_G/I_G^2\longrightarrow G,
 \qquad \sum_g n_gu_g+I_G^2\longmapsto\sum_g n_gg
\]
are mutually inverse homomorphisms. In degree zero the cyclic-chain
boundary is \(b(a\otimes b)=ab-ba\), so
\(HC_0^{\Z}(A_G)=A_G/[A_G,A_G]=A_G\). The homomorphism
\[
 \Theta_G:HC_0^{\Z}(A_G)\longrightarrow G,
 \qquad \sum_g n_gu_g\longmapsto\sum_g n_gg
\]
has kernel \(\Z u_0+I_G^2\) and fibre over \(g\) equal to
\((u_g-u_0)+\Z u_0+I_G^2\). Its restriction to \(I_G\) has
fibre \((u_g-u_0)+I_G^2\).
For every homomorphism \(f:G\to H\), the unital algebra map
\(f_*:A_G\to A_H\), \(u_g\mapsto u_{f(g)}\), preserves the
augmentation ideals and their squares and satisfies
\(\overline{f_*}\kappa_G=\kappa_H f\).
\end{lemma}
\begin{proof}
The displayed multiplication makes \(A_G\) a commutative unital ring.
An element of augmentation zero has the exact expression
\(\sum_g n_g(u_g-u_0)\). Consequently \(I_G^2\) is generated
as an additive group by the products
\[
 (u_g-u_0)(u_h-u_0)=u_{g+h}-u_g-u_h+u_0.
\]
The coefficient map \(\sum n_gu_g\mapsto\sum n_gg\) sends each
such product to \((g+h)-g-h+0=0\), and thus defines \(\theta_G\).
The same relation proves
\(\kappa_G(g+h)=\kappa_G(g)+\kappa_G(h)\).
Direct substitution gives \(\theta_G\kappa_G(g)=g\).
For \(w=\sum n_gu_g\in I_G\), additivity of \(\kappa_G\) gives
\[
 \kappa_G\theta_G(w+I_G^2)
 =\sum_g n_g(u_g-u_0)+I_G^2=w+I_G^2,
\]
where the final equality uses precisely \(\sum n_g=0\).
This proves the inverse and the restricted fibres. For arbitrary \(w\),
write \(w=\epsilon_G(w)u_0+(w-\epsilon_G(w)u_0)\). The second
summand lies in \(I_G\), while \(\Theta_G(u_0)=0\); hence the
kernel and every full fibre have the asserted form. Commutativity makes
the degree-zero boundary zero, proving the stated homology formula
directly. Finally the basis formula for \(f_*\) respects products,
unit, and augmentation. It sends \(I_G^2\) into \(I_H^2\),
and sends \(u_g-u_0\) to \(u_{f(g)}-u_0\). This proves naturality
on each coordinate, hence on every finite sum.
\end{proof}

\begin{theorem}[Cyclic characters with exact inverse and separation]
\label{thm:cyclic-character-bridge}
Let \(G\) be a finite abelian group and \(T=\Q/\Z\). Use the
cyclic cochain complex obtained by applying
\(\operatorname{Hom}_{\Z}(-,T)\) to the integral cyclic chain complex.
Its degree-zero cohomology consists of the additive traces
\(\tau:A_G\to T\), meaning \(\tau(ab)=\tau(ba)\).
Define the following subgroup of this degree-zero cohomology:
\[
 \mathcal T_G=\{\tau\in HC^0(A_G;T):
              \tau(u_0)=0,\quad \tau(I_G^2)=0\}.
\]
There is an isomorphism with specified inverse
\[
 \operatorname{Hom}(G,T)\longrightarrow\mathcal T_G,
 \quad \chi\longmapsto
 \left[\tau_\chi\!\left(\sum_g n_gu_g\right)=\sum_g n_g\chi(g)\right],
 \quad \chi(g)=\tau_\chi(u_g-u_0).
\]
Each fibre of this isomorphism is a singleton. These traces separate
the augmentation coordinates: \(g=0\) if and only if
\(\tau(u_g-u_0)=0\) for every \(\tau\in\mathcal T_G\).
If \(H\leq G\), every character of \(H\) extends to \(G\),
and the full fibre of restriction over a character \(\chi_H\) is
\(\chi_0+\operatorname{Hom}(G/H,T)\), where \(\chi_0\) is any
one extension and the second summand is pulled back along \(G\to G/H\).
\end{theorem}
\begin{proof}
The only cocycle condition in degree zero is vanishing on
\(ab-ba\), and there is no incoming cochain boundary in degree zero.
Since \(A_G\) is commutative every additive functional is a trace.
For a character \(\chi\), the formula for \(\tau_\chi\) is
additive, sends \(u_0\) to zero, and sends every generator
\(u_{g+h}-u_g-u_h+u_0\) of \(I_G^2\) to zero.
Conversely, for \(\tau\in\mathcal T_G\), set
\(\chi(g)=\tau(u_g-u_0)\). Applying \(\tau\) to that generator
proves \(\chi(g+h)=\chi(g)+\chi(h)\). Also \(\tau(u_0)=0\),
so additivity gives \(\tau(\sum n_gu_g)=\sum n_g\chi(g)\).
These substitutions prove both inverses and all singleton fibres.

Here is an explicit extension argument retaining all choices. Given
\(H\leq G\), \(\chi:H\to T\), and \(x\in G\), let \(m\)
be the least positive integer with \(mx\in H\); it exists because
\(G\) is finite. Choose a rational representative \(r\) of
\(\chi(mx)\). The solutions of \(m\beta=\chi(mx)\) in \(T\)
are exactly
\(\beta=(r+j)/m+\Z\), \(j=0,\ldots,m-1\): multiplying
verifies them, and subtracting two solutions shows their difference
is one of the \(m\) classes annihilated by \(m\).
For each such \(\beta\), define
\[
 \chi'(h+kx)=\chi(h)+k\beta\qquad(h\in H,\ k\in\Z).
\]
If \(h+kx=h'+k'x\), division of \(k-k'\) by \(m\) and
minimality of \(m\) show \(k-k'=\ell m\) for some integer
\(\ell\); then \(h-h'=-\ell mx\). The difference of the
two proposed values is \(-\ell\chi(mx)+\ell m\beta=0\).
Thus the formula is well defined and additive. It extends \(\chi\),
and every extension to \(H+\Z x\) is obtained uniquely by its
value \(\beta\) on \(x\). Adjoining an element outside the current
subgroup strictly increases its cardinality, so finitely many such
steps extend \(\chi\) to \(G\). The difference of two extensions
vanishes on \(H\), and factors uniquely through \(G/H\); adding
any such factor produces another extension. This proves the full fibre.

For nonzero \(g\), let its order be \(n>1\). The formula
\(kg\mapsto k/n+\Z\) defines a character on \(\langle g\rangle\):
two indices represent the same group element exactly when their
difference is divisible by \(n\). Extend this character by the
preceding construction. The resulting trace takes
\(u_g-u_0\) to \(1/n+\Z\ne0\). For \(g=0\), every such
evaluation is zero. This proves separation without discarding
the finite-order coordinates.
\end{proof}

\begin{theorem}[The original lattice obstruction in cyclic coordinates]
\label{thm:cyclic-lattice-gluing}
Keep the original \(q_1,\ldots,q_s,D,E,L,K_D,K_E,K_L\) and
the exact exponent box \(B\) of \eqref{eq:lattice-exact}. Set
\[
 G_L=\Z^s/K_L,\qquad
 G_0=(\Z^s/K_D)\oplus(\Z^s/K_E),\qquad
 O=\Z^s/(K_D+K_E).
\]
All three groups are finite. Under the coordinates of
Lemma~\ref{lem:integral-augmentation}, \eqref{eq:lattice-exact}
becomes the exact sequence
\[
 0\longrightarrow I_{G_L}/I_{G_L}^2
 \xrightarrow{\overline{i_*}} I_{G_0}/I_{G_0}^2
 \xrightarrow{\overline{\delta_*}} I_O/I_O^2
 \longrightarrow0,
\]
with the explicit generator maps
\[
 [u_{[e]_{K_L}}-u_0]\longmapsto
 [u_{([e]_{K_D},[e]_{K_E})}-u_0],\qquad
 [u_{([x]_{K_D},[y]_{K_E})}-u_0]\longmapsto
 [u_{[x-y]_{K_D+K_E}}-u_0].
\]
The corresponding cyclic-character maps form the exact sequence
\[
 0\longrightarrow\mathcal T_O
 \xrightarrow{\delta^*}\mathcal T_{G_0}
 \xrightarrow{i^*}\mathcal T_{G_L}\longrightarrow0,
 \qquad \delta^*(\tau)=\tau\circ\delta_*,\quad
 i^*(\tau)=\tau\circ i_*.
\]
For a specified local target pair
\(v=([x]_{K_D},[y]_{K_E})\), write
\(o=[x-y]_{K_D+K_E}\). Its integral gluing obstruction vanishes
if and only if every \(\tau\in\mathcal T_O\) has
\(\tau(u_o-u_0)=0\). If it vanishes, choose
\(x-y=k_D+k_E\) with \(k_D\in K_D,k_E\in K_E\), and set
\(e_0=x-k_D=y+k_E\). The complete original integer fibre and
its exact bounded part are, respectively,
\[
 e_0+K_L,\qquad B\cap(e_0+K_L).
\]
They use precisely the original coordinates and intervals; the bounded
part is enumerated by the fundamental-rectangle formula preceding
this construction. No nonemptiness of that intersection is assumed.
\end{theorem}
\begin{proof}
The map \([e]_{K_D}\mapsto\prod_j q_j^{e_j}\bmod D\) embeds
\(\Z^s/K_D\) in the finite unit group modulo \(D\): its kernel
is zero by the definition of \(K_D\). The same argument works
for \(E\) and \(L\). The group \(O\) is a quotient of
\(\Z^s/K_D\), hence is finite as well.
Lemma~\ref{lem:integral-augmentation} gives a bijective group
coordinate map \(\kappa\) at each term. Its naturality gives
the two displayed generator formulas and the commuting squares
with \(i\) and \(\delta\). Thus the middle kernel is exactly
\(\kappa_{G_0}(\ker\delta)=\kappa_{G_0}(\operatorname{im}i)\);
injectivity at the first term and surjectivity at the last follow
by applying the specified \(\theta\) inverses. Explicitly, the fibre
of \(\overline{\delta_*}\) over \(\kappa_O(o)\) is
\[
 \{\kappa_{G_0}(v_0+i(g)):g\in G_L\}
\]
for any \(v_0\) with \(\delta(v_0)=o\), and its parameter \(g\)
is unique because \(i\) is injective. Such a \(v_0\) exists:
if \(o=[t]_{K_D+K_E}\), take
\(v_0=([t]_{K_D},[0]_{K_E})\). A fibre of
\(\overline{i_*}\) over its image is a singleton, with inverse
\(\kappa_{G_L}i^{-1}\theta_{G_0}\); outside that image it is empty.
These are statements about the augmentation quotients. On them the
composite is zero, since \(u_0-u_0=0\), as the formulas show.

Under Theorem~\ref{thm:cyclic-character-bridge}, the dual maps are
\(\chi\mapsto\chi\circ\delta\) and
\(\psi\mapsto\psi\circ i\). Since \(\delta\) is surjective,
the first is injective. A character \(\psi\) on \(G_0\)
vanishes on \(i(G_L)=\ker\delta\) exactly when
\(\chi(o)=\psi(v)\), for any \(v\) with \(\delta(v)=o\),
defines a well-defined character on \(O\). This proves the middle
kernel and supplies the unique inverse on that image. Every character
on \(G_L\) extends from the subgroup \(i(G_L)\) to \(G_0\)
by the proved extension construction, so \(i^*\) is surjective.
Its full fibre over any specified trace is
\(\tau_0+\delta^*\mathcal T_O\), with \(\tau_0\) any one
extension; this follows from the computed kernel. The fibre of
\(\delta^*\) is a singleton on its image and empty elsewhere.

Character separation applied to \(O\) proves the criterion for
\(o=0\). If it holds, the displayed decomposition of \(x-y\)
gives \(e_0\equiv x\pmod{K_D}\) and
\(e_0\equiv y\pmod{K_E}\). A vector \(e\) has these same
two classes exactly when
\(e-e_0\in K_D\cap K_E=K_L\). This proves the entire fibre
and its intersection with \(B\), independently of the decomposition.
Each retained exponent vector reconstructs its original product
\(\prod_jq_j^{e_j}\); the original shell parameters
\(p,h,a,R_a,S_a\), when present, remain fixed throughout.
Every quotient used here has its specified inverse or complete fibre.
\end{proof}
